%\section*{Lernziele: Excel}
%\begin{todolist}
%\item Ich verstehe das Koordinatensystem von Excel-Zellen (z.B. $B4$ für eine Zelle in Kolonne 2, Zeile 4)
%\item Ich kann Zellen als Zahlen, Text, Währung, Datum und in weiteren Formaten darstellen.
%\item Ich kann einfach Formeln erstellen, beispielsweise um den Mittelwert, die Summe oder das Maximum einer Datenreihe zu bestimmen.
%\item Ich verstehe den Unterschied zwischen fixierten und nicht-fixierten Zellreferenzen (z.B. $B2$ vs. $\$B\$2$) und kann fixierte Zellreferenzen verwenden, um Formeln zu kopieren.
%\item Ich kann Datenreihen als Punkt-, Kreis- oder Liniendiagramm darstellen und kann, je nach Kontext, passende, weitere Diagrammelemente hinzufügen, wie beispielsweise eine Legende, X-Achsen-Titel, Y-Achsen-Titel, Grafik-Titel oder Führungslinien. Ich kann zudem den Bereich er X- und Y-Achse manuell anpassen.
%\end{todolist}

\section*{Lernziele: Aus Daten Lernen}
\begin{todolist}
\item Ich kann den Unterschied zwischen Kausalität und Korrelation anhand von konkreten Beispielen erläutern.
\item Ich kann das gemeinsame Verhältnis zwischen zwei Variablen A und B inhaltlich korrekt bestimmen, indem ich unterscheide folgende Fälle unterscheide:
\begin{todolist}
\item Direkte Kausalität ($A\rightarrow B$)
\item Indirekte Kausalität ($A\rightarrow C\rightarrow B$)
\item Gemeinsamer Grund ($C\rightarrow A$ und $C\rightarrow B$)
\item Zyklische Beeinflussung ($A\rightarrow B \rightarrow A$)
\item Koinzidenz (= zufällige Korrelation)
\end{todolist}
\item Ich weiss, was positive bzw. negative Korrelation ist und kann Beispiele für beide Fälle nennen.
\item Ich kann anhand von Punktdiagrammen erkennen, ob zwei Grössen positiv,
negativ oder gar nicht korrelieren (sofern dies einigermassen deutlich zu
erkennen ist).
\item Ich verstehe den Ausdruck $Z=x_1'\cdot y_1'+x_2'\cdot y_2'+...+x_n'\cdot y_n'$.
\item Ich kann den Unterschied zwischen $Z$ und dem Korrelationskoeffizienten $r$ erklären.
\item Ich kann den Korrelationskoeffizienten in einfachen Beispielen selbst berechnen
(mit TR).
\item Ich weiss, dass der Korrelationskoeffizient einheitenlos ist.
\item Ich weiss, dass der Korrelationskoeffizient immer zwischen -1 und +1 liegt.
\item Ich kann erklären, was mit einem \textit{Klassifikationsverfahren} gemeint ist.
\item Ich kann das \textit{K nearest neighbors}-Verfahren nutzen, um einen neuen Datenpunkt einer existierenden Klasse zuzuordnen.
\item Ich kann anhand einem konkreten Beispiel beschreiben, was passiert, wenn bei \textit{K nearest neighbors} ein grösserer oder kleinerer Wert für $K$ verwendet wird.
\item Ich kann den Unterschied zwischen \textit{K nearest neighbors} und einem Punkt-, bzw. Linien-Klassifikator anhand eines konkreten Beispiels erklären und kann den passerenden Klassifikator für die vorliegende Situation wählen.
\item Ich kann auf einer Grafik bestimmen, was Trainings- und was Testfehler sind.
\item Ich kann für ein gegebenes Datenset mit mehreren Attributen bestimmen, ob ein Punkt- oder Linienklassifikator die Daten eindeutig separieren kann.
\end{todolist}